\section{Five organizing principles}\label{sec:organizing-principles}

The mixed holomorphic-topological local model on
\(\R^2_{\mathrm{top}}\times\widehat{\C^2}_{0,\mathrm{hol}}\) is forced
by five principles.  Each is a single mathematical sentence; each
forces a named theorem in the manuscript.  Together they fix the open
brane complex, the closed Hamiltonian BF complex, the formal-local
CE/PV identification, the cotangent coefficient module, the support
mechanism of the factorization algebra, the scalar anomaly cocycle,
and the form in which every frontier obstruction is named.  The
classical attributions -- Dirac, Witten, Grothendieck -- enter as the
mathematical instruments inside the principles, not as the principles
themselves.

\subsection{Principle 1 (Brane Measurement Microscope)}

\emph{A stack of \(N\) Dirac branes on the topological line is a
finite Chan--Paton microscope for the formal holomorphic symplectic
disk; its measured observables are the stable scalar-reduced traces
\(J(f)=\overline{\operatorname{Tr}f(\phi_1,\phi_2)}\), and these
exhaust the gauge-invariant content of the probe pair.}

The substitution \(z_i\rightsquigarrow\phi_i\) of normal coordinates
into a commuting pair of \(\mathfrak{gl}_N\)-valued matrices is the
only datum: the \(GL_N\)-symplectic potential
\(\operatorname{Tr}(\phi_1\,d\phi_2)\), the moment map
\(-\operatorname{Tr}(\epsilon[\phi_1,\phi_2])\), the Dirac constraint
\([\phi_1,\phi_2]=0\), and the Koszul differential
\(Q\psi=[\phi_1,\phi_2]\) are its forced consequences.  The closed
Hamiltonian Lie algebra \(\bar A=\C[z_1,z_2]/\C\cdot1\) acts on the
brane through nothing other than \(J\); cyclic invariance and
adjacent-swap \(Q\)-exactness collapse every trace word to the
symmetric symbol family
\(\operatorname{Tr}(\phi_1^k\phi_2^l)=J(z_1^kz_2^l)\).  This forces
the matrix-microscope theorem
(Lemma~\ref{lem:probe-trace-microscope}) and, after formal Darboux
normalization at any brane point of any connected
holomorphic-symplectic surface, the measurement-universality theorem
of \S\ref{ssec:twistedM-universality}: every admissible protected
Hamiltonian sector reduces to the same brane microscope, and divergence
is read coordinate by coordinate in the four-class residual
\(\operatorname{Ob}^{\mathrm{univ}}_{p,X}\).

\subsection{Principle 2 (Three Localities)}

\emph{The factorization-algebra structure on the brane line is forced
by three independent local-current identities,
\[
  \delta(t-s)\quad\text{(support)},\qquad
  L_v=[Q,\iota_v]\quad\text{(topological)},\qquad
  L_{\partial/\partial\bar z_i}
   =[Q,\iota_{\partial/\partial\bar z_i}]
   \quad\text{(formal-holomorphic)},
\]
and none of the three implies the others.}

Support locality identifies \(P_0\)-brackets of disjoint smeared
Hamiltonian currents with zero; topological locality makes the brane
factorization algebra locally constant in \(\R^2_{\mathrm{top}}\);
formal-holomorphic locality concentrates closed-side cohomology on
holomorphic jets at the brane point and pairs cotangent modes with
Taylor Hamiltonians by residue.  Trace observables exhibit all three
identities; shifted-cotangent currents exhibit support and topological
locality but \emph{not} formal-holomorphic locality, and the gap is
precisely the cotangent centrality homotopy that no single identity
supplies.  This forces the kinematic locality theorem
(Proposition~\ref{prop:three-localities}) and the factorization-shape
determination of \(\mathcal F^{\mathrm{mix}}_{\mathrm{Ham}}\): the
unique factorization algebra compatible with the three identities has
dg Lie input
\(\Omega_c^\bullet(U)\widehat\otimes
\Omega_c^{0,\bullet}(B)\widehat\otimes
(\mathfrak h\ltimes\mathfrak h^\vee_{\mathrm{cont}}[1])\).

\subsection{Principle 3 (Defect Polarization)}

\emph{The brane defect line carries a canonical polarization:
configuration observables are polynomial Taylor classes
\(O_f=\overline{\operatorname{Tr}f(\phi_1,\phi_2)}\) on the
matrix-valued normal coordinate, while cotangent defect momenta are
distributional residues \(\rho_{a,b}=z_1^{-a-1}z_2^{-b-1}\,dz_1dz_2\)
in the Grothendieck--Matlis principal-part module
\(\mathcal P=\bigoplus_{a+b>0}\C\,\rho_{a,b}\).}

The continuous dual of the reduced Hamiltonian disk algebra is not
another copy of polynomials but the residue annihilator of the
constant line in the top local-cohomology module
\(H^2_{\mathfrak m}(\C[[z_1,z_2]])\,dz_1dz_2\); the residue pairing
makes the identification canonical, and the residue-coadjoint formula
\(z_1^p z_2^q\cdot\rho_{a,b}=-(pb-qa+p-q)\,\rho_{a-p+1,b-q+1}\) is its
closed-form fingerprint.  The polynomial PBW one-\(\psi\) descendants
\(\widetilde\Psi_{a,b}\) carry the same labels but a strictly different
\(\mathfrak h\)-module structure: linear Hamiltonians act locally
nilpotently on the polynomial side and never locally nilpotently on
\(\mathcal P\), so no relabelling of bases interpolates -- the PBW
classes are the special-fibre shadow, not the deformation-level
cotangent module.  This forces the defect polarization theorem
(Theorem~\ref{thm:app-defect-polarization}), the residue-coadjoint
uniqueness theorem
(Theorem~\ref{thm:principal-part-coadjoint-uniqueness}), and the type
discipline of the CE/PV dictionary
(Theorem~\ref{thm:reduced-ce-pv-central-operation}): boundary
Hamiltonian observables are the image of the cotangent coordinate
\(u\), not the Hamiltonian ghost coordinate \(c\).

\subsection{Principle 4 (Scalar Anomaly Transgression)}

\emph{The constant Hamiltonian generates the zero Hamiltonian vector
field closed-side and survives open-side as the brane \(U(1)\)
center of mass; the same scalar cocycle is detected algebraically as
\([\bar c]\in H^2_{\mathrm{Lie}}(\bar A,\C)\), Lie-theoretically as the
central extension of \(\bar A\), in the trace algebra as
\(\{\operatorname{Tr}\phi_1,\operatorname{Tr}\phi_2\}=N\), and after
Moyal quantization as the Capelli/QME shift \(\hbar N[\bar c]\).}

The four detections are not four cocycles -- they are one cocycle
transgressed across four faces of the same scalar-reduction step.
Removing constants from \(\C[[z_1,z_2]]\) to obtain \(\mathfrak h\)
projects out the kernel of the Hamiltonian vector field map
\(f\mapsto X_f\); on the brane the same kernel reappears as the
center-of-mass mode, so the closed-open mismatch is forced and equal
to the transgressed class.  Algebraic deformation quantization of the
disk by the Moyal star product converts the same scalar into the
order-\(\hbar\) Capelli term, with \(\hbar N\) the only admissible
coefficient; the first scalar-reduced \(\mathfrak{gl}_N\) QME
obstruction is therefore \(W_1^{\mathrm{sc,ord}}=\hbar N\,[\bar c]\),
and any cancellation must run through the unreduced central-character
branch, the opposite-extension branch, or the balanced-supertrace
branch.  This forces the scalar anomaly transgression theorems
(Theorems~\ref{thm:u1-center-anomaly},
\ref{thm:u1-center-anomaly-open},
\ref{thm:quantum-classical-anomaly}) and the cancellation theorems of
the unreduced QME appendix.

\subsection{Principle 5 (Supremum Form of Obstructions)}

\emph{Every frontier obstruction in this manuscript is stated as a
named cohomology class or a named contracting-homotopy obligation,
never as a hedge; ``conditional'', ``external'', and ``expected'' are
not stopping categories.}

Four frontier surfaces share this discipline.  The unreduced
\(P_0\)-factorization-centre lift is the comparison cone class
\([\Phi_{\mathrm{cl}}]\) of Problem~\ref{prob:formal-factorization-center},
with named boundary representatives, centrality homotopies, coefficient
kernel, and brane--defect QME solution as its construction data.  The
non-scalar Costello QME on a brane interval is the single class
\([\operatorname{Ob}^{\mathrm{red}}_{w,\partial,\hbar}]
\in H^1(\ker\widehat\sigma_{\mathrm{sc}},d_{\mathrm{QME}})\) together
with its \(\varprojlim^1H^0\)-primitive-compatibility pair, and its
first exposed larger row is \(\theta_3\) with covector
\(\ell_{\theta_3}(\mathsf E_{\theta_3,M})=1\) -- a named
\(C_{\theta_3}\) is evidence only when paired with a CE ancestor, a
scalar-zero Costello local counterterm, or a complete companion-face
table.  The strict native all-window Bochner--Martinelli current
transfer is the obstruction vector
\(\operatorname{Ob}^{\Pi}_{\mathrm{BM}}\), and the residue-coadjoint
formula on \(z_1^{N+2}\rho_{N+1,0}\) computes its first nonvanishing
component as \((N+1)(N+2)\rho_{0,1}\).  The all-bidegree decorated PBW
Stokes identity is the family of single classes
\(\Omega^{\mathrm{rad}}_{a,b}\in\operatorname{coker}B_{a,b}\),
equivalently the cokernel of the radial PBW boundary operator at each
bidegree, and a finite normal-form certificate is not inflated into
all-bidegree closure.  In every frontier surface the admissible repair
is therefore the smallest enlargement of the mathematical habitat in
which the intended mechanism is a theorem, or else the proved
obstruction theorem naming the exact class that prevents it.

\subsection{Forcing the local theorem}

The five principles have disjoint roles.  Principle~1 selects the open
BV complex and the stable trace observables.  Principle~2 selects the
factorization habitat and forces \(\mathcal F^{\mathrm{mix}}_{\mathrm{Ham}}\)
as its unique solution.  Principle~3 selects the cotangent coefficient
module and the operadic type discipline of the CE/PV dictionary.
Principle~4 inscribes the scalar cocycle as the only mismatch the
closed-open reduction can carry.  Principle~5 sets the form in which
every remaining frontier statement is named.  Their combined force
splits the local theorem into independent components -- the Dirac probe
theorem, the matrix-microscope and one-\(\psi\) calculation, the CE/PV
central-operation theorem, the Matlis principal-part theorem, the
brane-line factorization-current theorem, the reduced principal-part
defect-current theorem, the degree-zero Moyal/Weyl quantum theorem,
and the scalar \(U(1)\) anomaly theorem -- and isolates the four
frontier components as named obstruction surfaces governed by the
obstruction complex of \S\ref{sec:open-obligations}.
